\section{System Model}\label{sec:system_model}

\subsection{Terminal Position}
Let $\mathbf{x} = [e, n]^{\top}$ denote the unknown terminal easting-northing (EN)
offset in meters from a known reference point within a bounded region of interest
$\mathcal{R} \subset \mathbb{R}^2$.  The terminal altitude is assumed known or
fixed for this analysis.

\subsection{Satellite State from TLE/SGP4}
A ground station receives a TLE set for the visible LEO satellite.  The SGP4
propagator \cite{todo_sgp4} returns the satellite position
$\mathbf{s}(t) \in \mathbb{R}^3$ (km, ECEF) and velocity
$\mathbf{v}_s(t) \in \mathbb{R}^3$ (km/s, ECEF) at UTC time $t$.

\subsection{Geometric Range and Range-Rate}
The geometric range from terminal at ECEF position $\mathbf{p}_t$ to satellite is
\begin{equation}
  \rho(t; \mathbf{x}) = \bigl\|\mathbf{s}(t) - \mathbf{p}_t(\mathbf{x})\bigr\|,
\end{equation}
where $\mathbf{p}_t(\mathbf{x})$ is the terminal ECEF position corresponding to EN
offset $\mathbf{x}$ via a local ENU basis.
The range-rate is
\begin{equation}
  \dot{\rho}(t; \mathbf{x}) = \frac{(\mathbf{s}(t)-\mathbf{p}_t)\cdot\mathbf{v}_s(t)}
                                 {\rho(t; \mathbf{x})}.
\end{equation}

\subsection{Doppler Observation}
The one-way Doppler shift at carrier frequency $f_c$ is
\begin{equation}
  f_D(t; \mathbf{x}) = -\frac{f_c}{c}\,\dot{\rho}(t; \mathbf{x}),
\end{equation}
where $c = 299\,792.458$~km/s is the speed of light.

The terminal's low-quality oscillator contaminates the observed frequency with
a constant CFO $b_0$ (Hz), a linear drift term $b_1$ (Hz/s), and additive noise
$\epsilon_i \sim \mathcal{N}(0,\sigma_f^2)$.  At discrete observation time
$t_i = t_0 + i\Delta t_{\text{pkt}}$ the received frequency is
\begin{equation}
  y_f(t_i) = f_D(t_i; \mathbf{x}) + b_0 + b_1\,t_i + \epsilon_i .
  \label{eq:freq_obs}
\end{equation}

\subsection{Optional Delay Observation}
When a packet arrival timestamp is recorded, the propagation delay
$\tau(t) = \rho(t; \mathbf{x})/c$ is observable up to an unknown time offset
$\Delta t$ and noise $\eta_i \sim \mathcal{N}(0,\sigma_\tau^2)$:
\begin{equation}
  y_\tau(t_i) = t_i + \Delta t + \frac{\rho(t_i; \mathbf{x})}{c} + \eta_i .
  \label{eq:delay_obs}
\end{equation}
In low-cost receivers, timestamp resolution may be limited; the estimator
accounts for this via the noise std $\sigma_\tau$.

\subsection{Nuisance Parameter Summary}
\begin{table}[h]
  \caption{Nuisance and State Parameters}
  \label{tab:params}
  \centering
  \begin{tabular}{cll}
    \hline
    Symbol & Description & Units \\ \hline
    $\Delta t$ & Timing offset & s \\
    $b_0$      & Constant CFO    & Hz \\
    $b_1$      & Linear drift    & Hz/s \\
    $\epsilon$ & Freq.\ noise    & Hz \\
    $\eta$     & Delay noise     & s \\ \hline
  \end{tabular}
\end{table}
The unknown vector $\boldsymbol{\theta} = [\Delta t, b_0, b_1]^{\top}$ is treated
as nuisance to be marginalized or fitted jointly with position.